\section{Exploratory concentration and constrained transport}\label{app:formation}
A distinct formation hypothesis is that an initially disordered scalar
field amplifies a band of spatial wavelengths and subsequently fragments.
Condensate fragmentation and a scale selected by growing perturbations
are established questions in other Q-ball models~\cite{KasuyaKawasaki2000}.
They do not establish formation in our initial-data family or select the
BIC frequencies. A future comparison should vary narrow-band and broadband
initial spectra at matched total energy and charge, measure growth and
outgoing flux, and distinguish transient concentrations from persistent
localized charged objects. A connected high-density region in a periodic
box should also be checked for closed paths with nonzero winding around
the box. A region with such a path is not an isolated density island at
that threshold. Absence of winding alone would still not establish vacuum
isolation, stationarity or stability; threshold and resolution dependence
must be checked. Total charge conservation forbids production
of net charge from neutral data; opposite charges or a compensating
background would have to be accounted for. No universal integer multiple
of a fundamental wavelength is derived here.

Even persistence alone can be biased by the preparation. For the potential
in Eq.~\eqref{eq:action}, put $S=|\phi|^2$, $\Pi=\partial_t\phi$,
$Q=2\operatorname{Im}\int\phi^*\Pi\,\dd^3x$ and $q=|Q|>0$.
For finite total energy $E=\int(|\Pi|^2+|\nabla\phi|^2+\mathcal U(S))\dd^3x$,
the elementary identity
\[
 E-q=\int\!\left[|\Pi-\ii\operatorname{sgn}(Q)\phi|^2
       +|\nabla\phi|^2-S^2+\tfrac12S^3\right]\dd^3x
\]
implies $\int S^2\dd^3x\ge q-E$ when $E<q$.
Since $\mathcal U(S)\ge S/2$, it also gives
$\operatorname{ess\,sup} S\ge(q-E)/(2E)>0$.
Thus conserved whole-system energy and charge already preclude uniform
dilution to arbitrarily small amplitude for such initial data. This elementary
bound is not a formation result: it specifies neither a
fixed spatial core nor a single stationary Q-ball. Charge transport,
fragmentation and continued oscillation remain possible. A formation test
must therefore resolve localization, profile and phase dynamics rather
than count survival alone. The bound concerns the closed system, not an
unbalanced subvolume or a dissipative boundary model.

\paragraph{A separate finite-time radial relaxation study.}
A subsequent project study of the same potential evolved four prescribed
spherically symmetric Gaussian initial states to $T=1000$, with an
absorbing exterior. Its predeclared diagnostic compared the magnitude of
the charge inside $r<60$ with the stationary-family charge at the
frequency estimated from the central phase, allowing a relative
difference below $10\%$. This tests charge--frequency consistency,
not equality of the evolving field to a stationary solution.
The matched-width state passed at $T=500$ and the broader state at
$T=1000$. The narrower state failed and retained about $2\%$ of its
initial charge inside the measurement region. A larger state remained
strongly pulsating and also failed this diagnostic. Consequently, the
combined formation criterion fixed before the runs was not met.
The successful initial states already had whole-system $E/|Q|<1$;
they do not demonstrate capture of initially unbound radiation.
Two jointly refined spatial and temporal resolutions gave the same
classifications, but absorber reflection was not measured. These reported
finite-time radial observations neither establish exact relaxation nor
test nonspherical stability; a charge-retention plateau alone is not a
binding proof. The original plan and result are identified in the
reproducibility inventory. They have received a document-scope review,
not an independent reproduction of the trajectories.

% Insert after the formation/non-dilution discussion, before Remaining tasks.
% Exploratory numerical diagnostic; distinct from the certified BIC background.
\paragraph{One angular direction during chirped concentration.}
To examine a freedom excluded by radial evolution, we evolved one prescribed
quadrupolar tangent on a time-dependent inward chirp. This is an exploratory
formation diagnostic on a different background from the stationary BIC
studied above. In the units of Eq.~\eqref{eq:action}, the radial initial data are
$F(r)=(3/5)^{3/2}\exp(-r^2/50)\exp(\ii\beta r^2)$,
$\phi_b(r,0)=F(r)$ and $\partial_t\phi_b(r,0)=0.8\ii F(r)$,
with $\beta=0.015$. Their whole-space analytic values are
$Q\simeq240.55177$ and $E/Q\simeq1.038669965>1$; hence the preceding
$E<|Q|$ non-dilution argument does not force a remnant. The calculation
uses a radial Dirichlet box $R=80$ and ends at $T=32$.

For $\delta\phi=\eta(r,t)Y_{20}$, with real unit-normalized $Y_{20}$,
put $w=r\eta$ and $S_b=|\phi_b|^2$. The laboratory-frame tangent equation is
\begin{equation}
 w_{tt}=w_{rr}-\frac{6w}{r^2}
 -(1-4S_b+\tfrac92 S_b^2)w
 -(-2+3S_b)\phi_b^2\overline w .
 \label{eq:angular-chirp-tangent}
\end{equation}
We fix $w(r,0)=r^3(-1/25+2\ii\beta)F(r)$ and
$w_t(r,0)=0.8\ii w(r,0)$, the regular $O(r^3)$ quadrupolar deformation.
The nonlinear radial background is advanced simultaneously, without angular
backreaction. Its stored diagnostics reproduce the earlier chirp evolution;
this replay is a control, not independent formation evidence.

Centered second differences and RK4 are compared at the fixed choices
\[
(h,\Delta t)=(0.1,0.02),(0.05,0.02),(0.1,0.01).
\]
Origin polynomial, independent free-evolution and conjugate-coupling controls
precede the calculation. Our positive reduced diagnostic is
\begin{equation}
 N_h^2=h\sum_{j=1}^{J-1}
 \left[|w_{t,j}|^2+\left(1+\frac6{r_j^2}\right)|w_j|^2\right]
 +\frac1h\sum_{j=0}^{J-1}|w_{j+1}-w_j|^2,
 \qquad r_J=80 .
 \label{eq:angular-chirp-norm}
\end{equation}
It is not a conserved perturbation energy. Positive core, exterior and
collar contributions use fixed intervals $[0,6]$, $[6,76]$ and $[76,80]$,
with shared-node half weights and complete edge contributions.

On the finer spatial grid, the largest \emph{sampled} $N_h(t)/N_h(0)$ is
$1.356829$ at $t=21.6$; its value at $T=32$ is $1.139233$.
Samples are separated by $0.8$, so this is not a continuous-time supremum.
The deformation concentrates in the core and subsequently redistributes
outward (Fig.~\ref{fig:angular-chirp-l2}). Common-grid endpoint differences,
measured in Eq.~\eqref{eq:angular-chirp-norm} and divided by the base-grid
initial norm, are $0.0062849$ for spatial refinement and
$4.0021\times10^{-7}$ for timestep halving. These are sensitivity measures,
not certified error bounds. The sampled collar contribution satisfies
$N_{\mathrm{collar}}^2(t)/N_h^2(0)\le10^{-6}$, the prescribed
contamination threshold. This single direction shows transient amplification,
not a stability classification: the background remains nonstationary,
and other radial shapes, angular sectors, finite amplitudes and later times
remain untested.

\begin{figure}[htbp]
\centering
\includegraphics[width=.98\linewidth]{angular-chirp-l2.pdf}
\caption{Angular response of the inward chirp. Top: the actual signed linear
 density variation $\delta S=2\operatorname{Re}(\overline{\phi_b}\eta)Y_{20}$
 from fine-grid snapshots at $t=0,8,16,24,32$, in the fixed $xz$ window
 $[-12,12]^2$. Radial nodal values are interpolated for rendering; the common
 symmetric colour scale is per unit tangent amplitude. Bottom: global norm
 for all three discretizations, and fine-grid core/exterior squared-norm
 contributions divided by the \emph{global initial} $N_h^2(0)$.
 These are linear tangent data, not a finite-amplitude nonlinear
 three-dimensional evolution or a physical energy partition.}
\label{fig:angular-chirp-l2}
\end{figure}

\paragraph{Leading angular feedback on core-charge transport.}
A separate second-variation calculation follows
$\phi_\varepsilon=B+\varepsilon\eta Y_{20}+\varepsilon^2\Xi+O(\varepsilon^3)$
and evolves the forced spherical mean $\zeta=\langle\Xi\rangle_\Omega$.
Here $\Xi$ is half the second derivative, and $\zeta$ is an angular mean,
not the coefficient of normalized $Y_{00}$. Only this mean is needed
alongside the first tangent for the angle-integrated charge through
second order. The initial family is the volume-preserving ellipsoidal
deformation $r_\varepsilon^2=e^{-2c\varepsilon}(x^2+y^2)+e^{4c\varepsilon}z^2$,
$c=\sqrt{5/(16\pi)}$, with an explicitly differentiated finite-grid
normalization fixing total charge. For $Q_2$, the coefficient of
$\varepsilon^2$, define
\[
 b=\frac{Q_{2,\mathrm{core}}(32)-Q_{2,\mathrm{core}}(0)}{Q_0}
\]
in the same fixed core $r<6$, box $R=80$ and time window.
The three $(h,\Delta t)=(0.1,0.02),(0.05,0.02),(0.1,0.01)$ calculations
give $b=0.08737135,\ 0.08736583,\ 0.08737135$, respectively;
the spatially refined value exceeds the predeclared heuristic guard
$2.864\times10^{-5}$. This guard combines measured spatial, temporal and
charge-balance sensitivities with a fixed floor; it is not a certified
error bound or confidence interval. Stage-integrated physical charge flux
checks the core-charge difference, rather than defining it from the endpoint.

The interpretation requires the initial redistribution
(Fig.~\ref{fig:angular-second-order}). On the finer spatial grid,
$Q_{2,\mathrm{core}}(0)/Q_0=-0.10587389$, whereas
$Q_{2,\mathrm{core}}(32)/Q_0=-0.01850805$.
The radial reference itself has positive net core inflow over this window:
its core fraction increases from $0.58949606$ to $0.65267983$.
The positive second-order net-transport coefficient recovers much of an
initial core-charge deficit but leaves the final correction negative;
it does not show more final charge retained than in the radial reference.
The sampled response is nonmonotonic, with a transient positive core
correction before the negative final value; the decision uses the fixed
endpoint $T=32$, not a selected intermediate maximum.
The preparation fixes charge, not energy:
$E(\varepsilon,0)=E_0+\varepsilon^2 E_2+O(\varepsilon^3)$ with
$E_2/E_0=0.0224455$ on that grid. These are second-order Taylor coefficients at zero deformation,
not a finite-amplitude nonlinear three-dimensional evolution, a binding
comparison at equal energy, or evidence of stationary Q-ball formation.
As a subsequent analytical interpretation, the untruncated continuum
family has $E_{\mathrm{grad}}(\varepsilon)=
E_{\mathrm{grad},0}[2e^{-2c\varepsilon}+e^{4c\varepsilon}]/3$,
while its kinetic and potential integrals remain unchanged. Thus its
positive $E_2=4c^2E_{\mathrm{grad},0}$ is a geometric preparation cost;
this continuum identity neither replaces the stored discrete $E_2$ nor
constitutes a new acceptance criterion or determines the transport sign.

\begin{figure}[htbp]
\centering
\includegraphics[width=.98\linewidth]{angular-second-order-transport.pdf}
\caption{Second-order angular response from stored samples. Left:
$Q_{2,\mathrm{core}}(t)/Q_0$ for all three discretizations, with the initial
fine-space correction marked separately. Right: the initial-subtracted
coefficient $b(t)$ and the negative, independently stage-integrated
outward charge flux divided by $Q_0$, on the fine-space grid.
These routes check the same charge balance, not separate
physical replications. Straight lines join 41 stored samples; no finite
$\varepsilon$ is reconstructed. The final core correction remains negative
despite positive net transport. The endpoint guard is heuristic, not a
confidence interval, certified error or time-uniform uncertainty band.}
\label{fig:angular-second-order}
\end{figure}

% Numerical result independently reviewed: RESULT-REVIEW.txt cee669eda28b331f143750ed01913dc63e985b13ace98badac8e0c3458a652d3.
\paragraph{Compensating the initial energy through second order.}
A further, separately specified family keeps the same first angular
variation but replaces the initial chirp by
$\beta(\varepsilon)=\beta_0+\beta_2\varepsilon^2$.
On each grid, $\beta_2=-E_2/(\partial_\beta E_0)$ is determined from the
preparation energy, without using the later transport. The normalized
family fixes total charge exactly and cancels the initial energy change
through order $\varepsilon^2$; it does not assert exactly equal energies
at finite deformation. Only the homogeneous radial chirp tangent is newly
evolved, and its contribution is combined with the stored angular
second variation. The earlier fixed-charge result remains unchanged.

On the spatially refined grid, $\beta_2=-0.0082897517$ and the normalized
net-transport derivative with respect to $\beta$ is $-11.0018872$.
Their product adds $0.0912029133$ to the previous transport coefficient:
$b_{QE}=0.1785687477$, compared with $b_Q=0.0873658344$.
The base and time-refined values are $0.1785384553$ and $0.1785384478$;
the predeclared heuristic guard for the spatially refined transport
coefficient is $1.525216\times10^{-4}$, not a certified error interval.
The initial core correction is unchanged, while the final core coefficient
$c_{QE}=Q_{2,\mathrm{core},QE}(32)/Q_0$ is now $+0.0726948606$,
compared with $-0.0185080526$ for the previous family.
This final coefficient is reported separately from the
net-transport decision; no additional endpoint-success gate was introduced.
The compensation also reduces the initial inward radial current through
its negative chirp correction. Relative to the radial reference, this
energy-compensated family changes geometry and radial phase together;
its transport response cannot isolate geometry alone.
These coefficients concern the fixed endpoint $T=32$ in the same finite
box, not a finite-amplitude deformed solution, stationary formation or binding.
Figure~\ref{fig:fixed-energy-transport} separates the core coefficients
from their initial-subtracted transport for the two preparations.

\begin{figure}[htbp]
\centering
\includegraphics[width=.98\linewidth]{fixed-energy-transport.pdf}
\caption{Stored spatially refined second-variation coefficients with and
without initial-energy compensation. Left: the core-charge correction;
right: its change from the common initial correction. Straight segments
join 41 stored samples, without fitting or a new evolution. The original
family fixes charge; the compensated family also fixes initial energy
through order $\varepsilon^2$ while changing the radial chirp. Neither
curve represents a finite deformation or a percentage increase. Endpoint
caps on the transport panel are heuristic guards, not confidence intervals,
core-coefficient error bars or time-uniform bounds. The sign change of the
final core coefficient is not a binding or stability result.}
\label{fig:fixed-energy-transport}
\end{figure}

% Stored-state local-clock diagnostic, independently reviewed; finite scope only.
\paragraph{A local two-observable time-shift test.}
The energy-compensated response need not describe a new stationary state:
even an increased endpoint core coefficient could reflect a shifted phase
of the transient concentration. A fixed stored-state check at $T=32$
therefore compares two radial observables,
$C=Q_{\mathrm{core}}/Q_0$ and $H=R_{\mathrm{RMS}}^2/6^2$, where $R_{\mathrm{RMS}}^2$
is the second radial moment weighted by $|\phi|^2$ and divided by its
spatial integral. Write their second-order coefficients as $C_2,H_2$.
Here $C_2(32)=c_{QE}$, not the initial-subtracted $b_{QE}$.
A common local shift $t\mapsto t+\varepsilon^2\tau$ would require
$(C_2,H_2)=\tau(\dot C_0,\dot H_0)$, hence
\[
 \mathcal D=\dot C_0 H_2-\dot H_0 C_2=0.
\]
The moment and its derivative were evaluated directly from the stored
Cauchy data, including the norm-denominator variation; no time-series fit
or additional field evolution was used. On the fine spatial grid,
$H_2=-1.0936132$ and $\mathcal D=0.0021220476$, exceeding the predeclared
heuristic sensitivity guard $6.820019\times10^{-5}$ in inverse model-time
units. Thus this pair is inconsistent with one common local clock shift
under that diagnostic rule. The charge-only shift coefficient is
$\tau_C=C_2/\dot C_0=-9.1439597$, after its separate denominator check:
it multiplies $\varepsilon^2$ and is not a measured finite time shift.
This is only a two-observable endpoint comparison, conditional on the
existing variational outputs; it establishes neither persistence of the
core enhancement nor binding, stationary formation or a full-field distance.
The sensitivity guard is not a certified error bound.

% Independently reviewed complete composite; two production assignments retained.
\paragraph{A finite ellipsoidal response at matched initial charge and energy.}
We additionally evolved finite, axisymmetric ellipsoidal preparations at
$\varepsilon=\pm0.02,\pm0.04$, with initial charge and full discrete energy
matched to each discretization's radial reference within the declared root
tolerance. Unlike the preceding second variation, this tests finite initial
data rather than energy compensation only through order $\varepsilon^2$.
At the fixed endpoint $T=32$ define
\[
 D(e)=\frac{C(+e)+C(-e)-2C(0)}{2e^2},\qquad C=Q_{\mathrm{core}}/Q_0.
\]
This is a symmetric final core response, not initial-subtracted net transport.
On the spatially refined grid,
$D(0.02)=0.0727746617$ and $D(0.04)=0.0730141978$, against the new
same-grid, same-timestep variational reference $c_{\mathrm{ref}}=0.0726948531$.
The prescribed heuristic envelope at $e=0.02$ is $1.2055721\times10^{-4}$;
the discrepancy $7.9808621\times10^{-5}$ plus this allowance meets the
predeclared one-percent agreement criterion. This envelope is not a certified
error bound or a confidence interval. The resolved-remainder criterion is
not met, so no measured Taylor order is claimed.

The separate core-fraction differences at $\varepsilon=+0.02,-0.02$ are
$2.8253326\times10^{-5}$ and $2.9966404\times10^{-5}$, respectively,
both positive in the stored endpoint data; the symmetric envelope is not
a separate certified bound for each sign. These are dimensionless charge-fraction
differences, not energy percentages; a positive symmetric coefficient alone
would not establish both signs. Figure~\ref{fig:fa-response} separates them.
The preparation changes geometry and radial chirp together. This finite-time
comparison establishes neither persistent enhancement, stationary formation,
binding nor general three-dimensional stability. The even-$\ell$, $m=0$
evolution excludes non-axisymmetric dynamics.

The complete twenty-trajectory matrix and separate variational reference
were assembled from two bounded production assignments: fourteen completed
trajectories and the reference were inherited unchanged, and six remaining
trajectories were completed separately. The first assignment retains its
\texttt{INCOMPLETE} status; the full matrix, not a favourable subset,
is used in the comparison.

\begin{figure}[htbp]
\centering
\includegraphics[width=.98\linewidth]{fa-response.pdf}
\caption{Finite symmetric endpoint coefficients for four discretizations
(left) and separate signed core-fraction differences on the spatially refined
grid (right). Caps are heuristic sensitivities, not confidence intervals.
Straight segments join stored time samples; no extrapolation or Taylor-order
fit is used. Initial redistribution is retained rather than subtracted.}
\label{fig:fa-response}
\end{figure}

\begin{figure}[htbp]
\centering
\includegraphics[width=.95\linewidth]{fa-fields-t0.pdf}
\caption{Stored initial amplitude (upper row) and relative phase (lower row)
for the radial and $\varepsilon=\pm0.04$ preparations. The meridional cut
uses $r\le12$ inside the computational domain $R=80$; the circle marks the
measurement radius $r=6$. The common amplitude scale is shared with
Fig.~\ref{fig:fa-fields-final}. Relative phase is
$\arg(\phi_\varepsilon\overline{\phi_0})$ in the common initial phase
convention, masked if either amplitude is below $10^{-4}$ of the shared
maximum. The centre $r<0.1$ is masked; grey is not zero amplitude or phase.}
\label{fig:fa-fields-initial}
\end{figure}

\begin{figure}[htbp]
\centering
\includegraphics[width=.95\linewidth]{fa-fields-t32.pdf}
\caption{The same stored-field reconstruction at $T=32$, with the common
scales and masks of Fig.~\ref{fig:fa-fields-initial}. Complex radial mode
coefficients are interpolated before evaluating the angular basis and phase.
These are axisymmetric amplitude/relative-phase cuts, not a general 3D
simulation, an absolute-phase observable, or a binding diagnostic.
Visually similar colours on the full $[-\pi,\pi]$ scale do not establish
exact phase equality.}
\label{fig:fa-fields-final}
\end{figure}

\paragraph{Phase rotation and the fixed-charge kinetic gap.}
A similar phase colour does not test rigid harmonic motion. We evaluated
necessary conditions on the six stored states with $\varepsilon=0,\pm0.04$
at $t=0,32$, using the complete $R=80$ box. For the reduced mode arrays
$U=r\phi$, $V=\partial_t U$, write
$\|X\|_h^2=4\pi h\sum_{j=1}^{J-1}\sum_\ell|X_{j\ell}|^2$,
$I=\|U\|_h^2$, $K=\|V\|_h^2$ and, for $I>0$,
$\omega_{\mathrm{fit}}=Q/(2I)$.
With $A_h(U)$ the actual discrete nonlinear acceleration, define
\[
 r_v=\frac{\|V-\ii\omega_{\mathrm{fit}}U\|_h^2}{K+\omega_{\mathrm{fit}}^2 I},
 \qquad
 r_a=\frac{\|A_h(U)+\omega_{\mathrm{fit}}^2 U\|_h^2}
                 {\|A_h(U)\|_h^2+\omega_{\mathrm{fit}}^4 I}.
\]
Rigid single-frequency motion requires both numerators to vanish. Initially
$r_v$ vanishes by construction, whereas $r_a\simeq0.0830$; the prescribed
phase velocity does not solve the stationary-profile equation. At $T=32$
the three states give $r_v=0.0282474$--$0.0282631$ and
$r_a=0.0586355$--$0.0588005$. These are squared-residual quotients, not
amplitude percentages or kinetic-energy fractions. This completed endpoint
postprocessing involved no further time integration. It provides neither a core-only classification,
a relaxation law, a certified continuum residual bound nor a general 3D
stability result.

The unnormalized velocity residual has a separate exact algebraic meaning.
In the positive canonical kinetic norm, minimizing velocities at fixed
instantaneous $U$ and total $Q$ gives $V_Q=\ii Q U/(2I)$ and
\[
 E_h(U,V)-E_{Q,h}[U]
 =J_v=\|V-\ii\omega_{\mathrm{fit}}U\|_h^2
 =K-\frac{Q^2}{4I},\qquad
 E_{Q,h}[U]=W_h(U)+\frac{Q^2}{4I},
\]
where $W_h$ contains the radial and angular gradient plus potential energy.
This is a comparison at fixed profile, not a physical projection or measured
relaxation: it generally changes regional charge and need not preserve
other conserved quantities. It does not minimize the spatial profile,
show how energy is radiated, or establish binding. In particular, $r_v$
normalizes $J_v$ by $K+Q^2/(4I)$, not by kinetic or total energy.

\clearpage
% Reviewed stored-data diagnostic; no new PDE evolution or scattering theorem.
% Stored-data diagnostic, not a new PDE evolution or scattering theorem.
\paragraph{Direction of the exterior field.}
For the same chirp, we analysed the stored radial Cauchy data
$u=r\phi_b$, $v=\partial_tu$ at $t=16,24,32$ on all three discretizations.
A fixed window $\chi$ rises from zero to one on $12<r<16$ as
$10s^3-15s^4+6s^5$, $s=(r-12)/4$; $(\chi u,\chi v)$ is zero-padded
onto a periodic analysis interval of length $L=320$. With Fourier series
normalized by $1/M$ for $M$ grid points, both temporal branches are retained:
$a_\sigma=(\widehat{\chi u}+\widehat{\chi v}/(\sigma\ii\Omega))/2$,
$\Omega=\sqrt{1+k^2}$, $\sigma=\pm1$.
The convention $\exp(\ii k r+\sigma\ii\Omega t)$ gives group velocity
$-\sigma k/\Omega$ and branch energy $8\pi L\Omega^2|a_\sigma|^2$.
The denominator $E_{\mathrm{win}}$ below is the \emph{total free spectral window
energy}, including the directionless $k=0$ and unresolved Nyquist bins;
it is neither the original nonlinear energy nor its escaped fraction.

The fine spatial grid gives the values below. Define
$\epsilon_{\mathrm{nl}}=\|\chi(-2S_b+\tfrac32 S_b^2)u\|_2/\|\chi u\|_2$
and $\epsilon_\chi=\|-2\chi'u_r-\chi''u\|_2/\|\chi u\|_2$,
using the radial discrete $L^2$ norm. These are operator-source quotients
in the mass-$1$ units of the model (dimensionally mass squared), not
relative energy errors.
\begin{center}
\begin{tabular}{rrrr}
\toprule
$t$ & $E_{\mathrm{out}}/E_{\mathrm{win}}$ & $\epsilon_{\mathrm{nl}}$ & $\epsilon_\chi$\\
\midrule
16 & 0.904702 & $4.55352\times10^{-5}$ & 0.593850\\
24 & 0.888886 & $1.14544\times10^{-4}$ & 0.828330\\
32 & 0.939306 & $1.58566\times10^{-3}$ & 0.573484\\
\bottomrule
\end{tabular}
\end{center}

All nine decompositions have $E_{\mathrm{out}}>E_{\mathrm{in}}+E_{\mathrm{Nyquist}}$
and $E_{\mathrm{out}}>E_{\mathrm{win}}/2$. The largest change in the outgoing
fraction under spatial refinement is $4.59\times10^{-4}$, and under
timestep halving $1.78\times10^{-8}$; these are sensitivity bounds over
the sampled cases, not certified errors. The exterior nonlinear source
is weak in these units, but the filter source is not: its quotient ranges
from $0.57$ to $0.83$. Zero extension also produces the separately monitored
boundary term $u_r(80^-)\delta(r-80)$, rather than another volume norm.
Thus the fixed-window data are predominantly outward-oriented; without
control of the accumulated source terms, this does not establish
asymptotically free radiation, stationary Q-ball formation or binding of
a remaining core.


\paragraph{A prescribed weak-potential extension.}
In an exploratory, prescribed instantaneous Klein--Gordon--Poisson closure,
matching candidates associated with the first two ladder labels persisted
at two weak couplings each and shifted to lower $\omega^2$. These are four
tested points, not a uniform robustness region. The observation remains
conditional on unresolved Poisson-residual and matching-map regularity
checks; it does not establish robustness in the Einstein--Klein--Gordon
system or nonlinear stability. Including the potential perturbation here
means response within the same instantaneous closure, not full general
relativity. The quantity $2|\Phi(0)|$ denotes twice the central potential
depth and is not identified with the compactness $2GM/R$.

